\documentclass[11pt,a4paper]{article}
\usepackage[T1]{fontenc}
\usepackage{lmodern}
\usepackage[margin=26mm]{geometry}
\usepackage{amsmath,amssymb,amsthm}
\usepackage{microtype}
\usepackage{enumitem}
\usepackage[hidelinks]{hyperref}
\setlength{\parskip}{3pt}
\setlength{\emergencystretch}{3em}
\newtheorem{theorem}{Theorem}[section]
\newtheorem{lemma}[theorem]{Lemma}
\newtheorem{corollary}[theorem]{Corollary}
\theoremstyle{remark}
\newtheorem{remark}[theorem]{Remark}
\newcommand{\Z}{\mathbb Z}
\newcommand{\N}{\mathbb N}
\newcommand{\Q}{\mathbb Q}
\newcommand{\F}{\mathbb F}
\newcommand{\ps}[1]{#1[[x]]}
\hypersetup{pdftitle={A396102: Congruences for an inverse-iterate family},pdfauthor={Dev Goyal}}
\title{A396102\\[3pt]Congruences for an inverse-iterate family}
\author{Dev Goyal}
\date{24 October 2026}
\begin{document}
\maketitle
\begin{center}
\fbox{\begin{minipage}{0.91\linewidth}\small
\textbf{AI contribution and review status.}
ChatGPT (OpenAI) generated the mathematical exposition, proof development,
Lean interfaces, and checks for this collection in response to Dev Goyal's
prompts. Dev Goyal is the sole human author. Independent human review is
pending. Original sequence authors, known mathematical results, and adapted
upstream code retain their separate credits. This document makes no claim
of publication, OEIS acceptance, or worldwide priority.
\end{minipage}}
\end{center}
\begin{abstract}
We construct the unique normalized integer solution of $A(A(A(x)))=(1+x)A(A(x))$ and prove that every positive-index coefficient is one modulo three. A parameter theorem determines every natural compositional iterate modulo the parameter, with an explicit source-equation and inverse bridge.
\end{abstract}
\section{A396102 and the inverse-iterate family}

For each positive integer $r$, consider the normalized series equation
\begin{equation}\label{eq:inverseoriginal}
F_r^{\circ r}(x)=(1+x)F_r^{\circ(r-1)}(x).
\end{equation}
The case $r=3$ is the original A396102 equation. We construct it through
an auxiliary series, rather than assuming that a solution exists.

\begin{lemma}\label{lem:inverseconstruct}
For every $r\ge0$ and every commutative coefficient ring, there is a
unique normalized series $R_r$ satisfying
\begin{equation}\label{eq:inverseaux}
R_r=x-R_r R_r^{\circ r}.
\end{equation}
If $r\ge1$, its compositional inverse $F_r$ is the unique normalized
solution of \eqref{eq:inverseoriginal}.
\end{lemma}

\begin{proof}
Define $T_r(G)=x-GG^{\circ r}$. It preserves normalization. If two
normalized series agree through degree $n-1$, their iterates do also.
The difference of the two products $GG^{\circ r}$ and $HH^{\circ r}$
has one factor divisible by $x^n$ and another divisible by $x$ in
each term. Thus
\[
G-H=O(x^n)\quad\Longrightarrow\quad
T_r(G)-T_r(H)=O(x^{n+1}).
\]
Starting at $x$ and iterating $T_r$, every coefficient stabilizes after
finitely many steps. The stabilized coefficients give a fixed point;
the same order-increase implication proves uniqueness. This is a
purely coefficientwise construction and works over arbitrary
commutative rings.

Every normalized series $U$ has a unique normalized compositional inverse.
Indeed, in $U(V)=x$ the coefficient of degree $n\ge2$ is $v_n$ plus
an expression in earlier coefficients of $V$, so it determines $v_n$
with multiplier one and no division. This constructs a right inverse
$V$. Construct a right inverse $W$ of $V$ in the same way; associativity
gives $U=U\circ(V\circ W)=(U\circ V)\circ W=W$, hence $V\circ U=x$
as well. The coefficient recursion proves uniqueness.

Take the genuine compositional inverse $F_r$ of $R_r$. Substitute
$F_r^{\circ r}(x)$ into \eqref{eq:inverseaux}. Inverse cancellation
gives
\[
F_r^{\circ(r-1)}=F_r^{\circ r}-xF_r^{\circ(r-1)},
\]
which is \eqref{eq:inverseoriginal}. Conversely, substitute
$R^{\circ r}$ into \eqref{eq:inverseoriginal}, where $R=F^{-1}$.
This gives $x=(1+R^{\circ r})R$, hence \eqref{eq:inverseaux}.
Uniqueness of $R_r$ now implies uniqueness of $F_r$.
\end{proof}

\begin{theorem}\label{thm:inversemod}
For every $r\ge1$,
\[
R_r(x)\equiv\frac{x}{1+x}\pmod r,\qquad
F_r(x)\equiv\frac{x}{1-x}\pmod r.
\]
Consequently, for $n\ge1$ and $k\ge0$,
\begin{equation}\label{eq:iteratemod}
[x^n]F_r^{\circ k}(x)\equiv k^{n-1}\pmod r.
\end{equation}
Here $0^0=1$, consistently with the degree-one coefficient of the
zeroth compositional iterate $x$.
\end{theorem}

\begin{proof}
Write $M_c(x)=x/(1-cx)$. Direct substitution gives
$M_a\circ M_b=M_{a+b}$ over any commutative ring. Therefore
$M_{-1}^{\circ r}=M_{-r}=x$ over $\Z/r\Z$. Moreover
\[
x-M_{-1}(x)x=M_{-1}(x),
\]
so $M_{-1}$ satisfies the auxiliary fixed-point equation modulo $r$.
The coefficientwise uniqueness in Lemma~\ref{lem:inverseconstruct}
identifies it with the reduction of $R_r$. Its actual inverse is
$M_1$, identifying the reduction of $F_r$. Finally,
$M_1^{\circ k}=M_k=x/(1-kx)$, whose degree-$n$ coefficient is
$k^{n-1}$ for $n\ge1$.
\end{proof}

For $r=3$ and $k=1$, this proves the A396102 conjecture
\[
[x^n]F_3(x)\equiv1\pmod3\qquad(n\ge1).
\]
The formal proof includes both directions of the correspondence with
the original triple-composition equation, not merely a congruence for
an auxiliary series. The formal parameterization uses $r+1$ when
stating the original equation so that its predecessor iterate is
always a natural number. The coefficient theorem itself also has a
well-defined $r=0$ version in \texttt{ZMod 0}; no such extension is
needed to settle the stated OEIS problem.

\section{Lean proof and exact scope}
The dedicated entry file is \path{lean/Entries/A396102.lean}. Its declarations lie in namespace \texttt{FinalCollection.A396102}. The principal checked endpoints are:
\begin{itemize}[nosep]
\item \texttt{source\_equation}: the exact triple-composition equation.
\item \texttt{source\_existence\_uniqueness}: existence and uniqueness of its normalized solution.
\item \texttt{coefficient\_congruence}: the source assertion for every $n\ge1$.
\item \texttt{iterate\_congruence}: every natural iterate modulo three.
\item \texttt{family\_source\_existence}: the entire family of original equations.
\item \texttt{family\_source\_equation}: the canonical family solves that equation.
\item \texttt{auxiliary\_source\_equation}: the fixed-point equation constructing the inverse.
\item \texttt{inverse\_right, inverse\_left}: both genuine inverse identities.
\item \texttt{family\_coefficient\_congruence}: every coefficient modulo the parameter.
\item \texttt{family\_iterate\_congruence}: the general iterate formula.
\end{itemize}
The full implementations are in \path{lean/More/InverseIterate.lean}.
The source offset is one. The formal original-family statement uses parameter $r+1$, while $A396102=F_3$. Thus $r=2$ in the original-family existence theorem gives the required triple equation. The congruence statements also allow the formal convention \texttt{ZMod 0}, but the mathematical modulus assertions here use positive parameters.

The companion is a separate, explicitly typed Lean entry file. Its proofs
invoke the preserved implementation modules; the complete dependency source
is included in the same archive. Shared lemmas are reused rather than
assumed. The entry file is not a self-contained replacement for Mathlib or
the supplied common library.

The project pins Lean 4.19.0 and Mathlib revision
\begin{center}\small\ttfamily
c44e0c8ee63ca166450922a373c7409c5d26b00b.
\end{center}
The collection checker recompiles the proof sources and audits the
theorems' transitive axiom dependencies. Only \texttt{propext},
\texttt{Classical.choice}, and \texttt{Quot.sound} are permitted.
The delivered verification reports record the actual successful runs.
Finite numerical checks are corroboration; the displayed conclusions are
unbounded theorems. An AI review of source correspondence does not replace
independent human review.
\begin{thebibliography}{9}
\raggedright
\bibitem{oeis} OEIS Foundation, entry A396102,
\url{https://oeis.org/A396102}. Source checked 24 October 2026.
The source sequence and selected conjectures are credited to Paul D. Hanna.
\bibitem{repo} D. Goyal, supplied reference repository,
\url{https://github.com/dvgyl/oeis-a396843-proof}.
This separate presentation belongs to the final collected proof package;
no earlier DOI is reassigned to it.
\end{thebibliography}
\end{document}
